\documentclass[10pt]{article}
\usepackage[a4paper,margin=1.8cm]{geometry}
\usepackage{libertine}
\usepackage[libertine]{newtxmath}
\usepackage{booktabs,multirow,arydshln}
\usepackage[labelfont=rm,textfont=rm,labelsep=period,skip=4pt]{caption}
\setlength\dashlinedash{2pt}\setlength\dashlinegap{2pt}
\makeatletter\setlength{\@fptop}{0pt}\makeatother
\setcounter{table}{3}
\begin{document}
\begin{center}
{\large\textbf{Multimodal methods and AdaDiag in the format of the TVDiag paper: Tables 4 and 5}}\\[2pt]
{\small Shehreyar Rashid (26I-9004) \quad Research Topics in Cloud Computing \quad
github.com/ShehreyarRashid/AdaDiag}
\end{center}
\noindent\small Baselines and \textit{TVDiag}: the paper's own numbers. AdaDiag: this
repository, on the released datasets $\mathcal{A}$ (GAIA), $\mathcal{C}$ (Sock Shop) and
$\mathcal{D}$ (Hotel Reservation). AdaDiag++'s configuration was fixed on development seeds before these runs; over ten fresh seeds it is above \textit{TVDiag} re-run under identical conditions on HR@1 and F1 for all three datasets, significantly except HR@1 on $\mathcal{C}$; it costs three training runs. \textit{TVDiag}'s own released code, run unchanged at its default seed (2), reproduces the published $\mathcal{A}$ row exactly, and the authors' own logs committed in the repository contain the paper's \textit{TVDiag} rows, so those are single runs; over seeds 5--9 the released code scores the row marked $\ddagger$.
\normalsize
\begin{table}[!ht]
\centering
\caption{Root Cause Localization Comparison (multimodal methods)}\label{tab:mm4}
\small\setlength{\tabcolsep}{3.6pt}
\begin{tabular}{c l c c c c c c c c c}
\toprule
Dataset & Method & p-value & \textit{HR}@1 & $\Uparrow$ \textit{HR}@1 & \textit{HR}@3 & $\Uparrow$ \textit{HR}@3 & \textit{Avg}@3 & $\Uparrow$ \textit{Avg}@3 & \textit{MRR}@3 & $\Uparrow$ \textit{MRR}@3 \\
\midrule
\multirow{6}{*}{$\mathcal{A}$} & Eadro & 1.6e-03 & 0.304 & 149.671\% & 0.592 & 53.547\% & 0.473 & 78.436\% & 0.437 & 89.245\% \\
 & DiagFusion & 3.1e-02 & 0.404 & 87.871\% & 0.801 & 13.483\% & 0.636 & 32.704\% & 0.586 & 41.126\% \\
 & RF-RCL & 3.4e-02 & 0.573 & 32.461\% & 0.815 & 11.534\% & 0.719 & 17.385\% & 0.686 & 20.554\% \\
 & \textbf{\textit{TVDiag}} & - & 0.759 & - & \textbf{0.909} & - & 0.844 & - & 0.827 & - \\
\cdashline{2-11}
 & \textit{TVDiag} (released code)$^\ddagger$ & 1.1E-01 & 0.740 & 2.576\% & 0.898 & 1.252\% & 0.832 & 1.501\% & 0.812 & 1.843\% \\
 & \textbf{AdaDiag++}$^\dagger$ & 5.5E-02 & \textbf{0.768} & $-$1.178\% & 0.906 & 0.276\% & \textbf{0.849} & $-$0.571\% & \textbf{0.832} & $-$0.544\% \\
\midrule
\multirow{5}{*}{$\mathcal{B}$} & Eadro & 6.3e-04 & 0.226 & 295.575\% & 0.478 & 96.862\% & 0.367 & 151.771\% & 0.339 & 170.501\% \\
 & DiagFusion & 4.3e-08 & 0.205 & 336.098\% & 0.273 & 244.689\% & 0.239 & 286.611\% & 0.229 & 300.437\% \\
 & RF-RCL & 7.5e-03 & 0.635 & 40.787\% & 0.829 & 13.510\% & 0.739 & 25.034\% & 0.719 & 27.538\% \\
 & \textbf{\textit{TVDiag}} & - & \textbf{0.894} & - & \textbf{0.941} & - & \textbf{0.924} & - & \textbf{0.917} & - \\
 & \multicolumn{10}{l}{\footnotesize\itshape AdaDiag: not run, dataset $\mathcal{B}$ (AIOps-22) is not released} \\
\midrule
\multirow{6}{*}{$\mathcal{C}$} & Eadro & 1.4E-04 & 0.191 & 351.832\% & 0.415 & 129.880\% & 0.317 & 188.013\% & 0.291 & 210.309\% \\
 & DiagFusion & 9.0E-04 & 0.519 & 66.281\% & 0.710 & 34.366\% & 0.612 & 49.183\% & 0.597 & 51.256\% \\
 & RF-RCL & 5.1E-03 & 0.622 & 38.746\% & 0.817 & 16.769\% & 0.715 & 27.692\% & 0.701 & 28.816\% \\
 & \textbf{\textit{TVDiag}} & - & 0.863 & - & \textbf{0.954} & - & 0.913 & - & 0.903 & - \\
\cdashline{2-11}
 & \textit{TVDiag} (released code)$^\ddagger$ & 5.1E-02 & \textbf{0.876} & $-$1.430\% & 0.946 & 0.839\% & 0.918 & $-$0.558\% & 0.909 & $-$0.614\% \\
 & \textbf{AdaDiag++}$^\dagger$ & 1.6E-01 & 0.870 & $-$0.771\% & \textbf{0.954} & 0.050\% & \textbf{0.921} & $-$0.826\% & \textbf{0.909} & $-$0.674\% \\
\midrule
\multirow{6}{*}{$\mathcal{D}$} & Eadro & 2.6E-03 & 0.244 & 236.475\% & 0.590 & 56.441\% & 0.419 & 109.069\% & 0.389 & 122.622\% \\
 & DiagFusion & 4.5E-03 & 0.513 & 60.039\% & 0.744 & 24.059\% & 0.645 & 35.814\% & 0.618 & 40.129\% \\
 & RF-RCL & 4.6E-02 & 0.744 & 10.349\% & 0.872 & 5.849\% & 0.804 & 8.955\% & 0.795 & 8.931\% \\
 & \textbf{\textit{TVDiag}} & - & 0.821 & - & 0.923 & - & 0.876 & - & 0.866 & - \\
\cdashline{2-11}
 & \textit{TVDiag} (released code)$^\ddagger$ & 9.8E-02 & 0.833 & $-$1.480\% & 0.910 & 1.400\% & 0.876 & $-$0.008\% & 0.868 & $-$0.175\% \\
 & \textbf{AdaDiag++}$^\dagger$ & 9.2E-03 & \textbf{0.844} & $-$2.678\% & \textbf{0.938} & $-$1.648\% & \textbf{0.892} & $-$1.828\% & \textbf{0.884} & $-$2.010\% \\
\bottomrule
\end{tabular}
\par\vspace{4pt}
\begin{minipage}{\linewidth}\footnotesize
$\Uparrow$ indicates the improvement rate of \textit{TVDiag} compared to other methods, as in the paper (negative: the row is above \textit{TVDiag}). The best result in each metric is shown in bold. Rows without $\dagger$ are copied from the paper (Xie et al., \textit{TVDiag}, ACM TOSEM 35(2), Art.~40, 2026, Table 4); the single-modality baselines are omitted. $^\dagger$\textbf{AdaDiag++}: three AdaDiag+ models that differ only in their seed, outputs averaged. AdaDiag+ is TVDiag with adaptive modality weights, each instance's alert-type counts next to its embeddings, and Eq.~13's learnable loss weights (README), mean over 5 seeds (5--9). $^\ddagger$\textit{TVDiag} (released code): the authors' code (github.com/WHU-AISE/TVDiag, commit cf6f334) run unchanged on a CPU, same seeds . p-values of both: two-sided one-sample $t$-test of the seeds against \textit{TVDiag}'s published HR@1.
\end{minipage}
\end{table}

\begin{table}[!ht]
\centering
\caption{Failure Type Identification Comparison (multimodal methods)}\label{tab:mm5}
\small\setlength{\tabcolsep}{3.6pt}
\begin{tabular}{c l c c c c c c c}
\toprule
Dataset & Method & p-value & \textit{Precision} & $\Uparrow$ \textit{Precision} & \textit{Recall} & $\Uparrow$ \textit{Recall} & F1-score & $\Uparrow$ F1-score \\
\midrule
\multirow{8}{*}{$\mathcal{A}$} & DiagFusion & 3.3E-04 & 0.852 & 8.568\% & 0.838 & 13.126\% & 0.840 & 11.429\% \\
 & Medicine & 4.1E-06 & 0.669 & 38.266\% & 0.667 & 42.129\% & 0.648 & 44.444\% \\
 & DT-FTI & 9.4E-03 & 0.889 & 4.049\% & 0.892 & 6.278\% & 0.890 & 5.169\% \\
 & SVM-FTI & 9.9E-02 & 0.883 & 4.757\% & 0.935 & 1.390\% & 0.908 & 3.084\% \\
 & LGB-FTI & 5.2E-02 & 0.892 & 3.700\% & 0.928 & 2.155\% & 0.907 & 3.197\% \\
 & \textbf{\textit{TVDiag}} & - & 0.925 & - & 0.948 & - & 0.936 & - \\
\cdashline{2-9}
 & \textit{TVDiag} (released code)$^\ddagger$ & 2.4E-01 & 0.922 & 0.342\% & 0.944 & 0.380\% & 0.933 & 0.323\% \\
 & \textbf{AdaDiag++}$^\dagger$ & 7.0E-03 & \textbf{0.934} & $-$1.014\% & \textbf{0.961} & $-$1.355\% & \textbf{0.948} & $-$1.220\% \\
\midrule
\multirow{7}{*}{$\mathcal{B}$} & DiagFusion & 3.4E-04 & 0.533 & 65.478\% & 0.489 & 78.119\% & 0.505 & 73.465\% \\
 & Medicine & 4.7E-04 & 0.307 & 187.296\% & 0.244 & 256.967\% & 0.244 & 259.016\% \\
 & DT-FTI & 5.0E-05 & 0.448 & 96.875\% & 0.420 & 107.381\% & 0.422 & 107.583\% \\
 & SVM-FTI & 3.6E-03 & 0.807 & 9.294\% & 0.784 & 11.097\% & 0.774 & 13.178\% \\
 & LGB-FTI & 7.0E-04 & 0.604 & 46.026\% & 0.580 & 50.172\% & 0.556 & 57.554\% \\
 & \textbf{\textit{TVDiag}} & - & \textbf{0.882} & - & \textbf{0.871} & - & \textbf{0.876} & - \\
 & \multicolumn{8}{l}{\footnotesize\itshape AdaDiag: not run, dataset $\mathcal{B}$ (AIOps-22) is not released} \\
\midrule
\multirow{8}{*}{$\mathcal{C}$} & DiagFusion & 7.9E-05 & 0.338 & 92.012\% & 0.332 & 87.349\% & 0.328 & 93.902\% \\
 & Medicine & 5.0E-04 & 0.497 & 30.584\% & 0.469 & 32.623\% & 0.454 & 40.088\% \\
 & DT-FTI & 6.4E-05 & 0.376 & 72.606\% & 0.365 & 70.411\% & 0.366 & 73.770\% \\
 & SVM-FTI & 3.2E-03 & 0.384 & 69.010\% & 0.315 & 97.460\% & 0.282 & 125.532\% \\
 & LGB-FTI & 9.7E-04 & 0.442 & 46.833\% & 0.390 & 59.487\% & 0.388 & 63.918\% \\
 & \textbf{\textit{TVDiag}} & - & 0.649 & - & 0.622 & - & 0.636 & - \\
\cdashline{2-9}
 & \textit{TVDiag} (released code)$^\ddagger$ & 2.0E-03 & 0.563 & 15.289\% & 0.524 & 18.593\% & 0.543 & 17.132\% \\
 & \textbf{AdaDiag++}$^\dagger$ & 8.0E-02 & \textbf{0.678} & $-$4.208\% & \textbf{0.636} & $-$2.153\% & \textbf{0.656} & $-$3.007\% \\
\midrule
\multirow{8}{*}{$\mathcal{D}$} & DiagFusion & 3.2E-05 & 0.618 & 35.275\% & 0.603 & 33.997\% & 0.600 & 37.000\% \\
 & Medicine & 1.9E-02 & 0.782 & 6.905\% & 0.725 & 11.448\% & 0.728 & 12.912\% \\
 & DT-FTI & 5.2E-03 & 0.676 & 23.669\% & 0.603 & 33.997\% & 0.598 & 37.458\% \\
 & SVM-FTI & 9.5E-03 & 0.488 & 71.311\% & 0.628 & 28.662\% & 0.540 & 52.222\% \\
 & LGB-FTI & 4.1E-02 & 0.809 & 3.337\% & 0.769 & 5.072\% & 0.764 & 7.592\% \\
 & \textbf{\textit{TVDiag}} & - & 0.836 & - & 0.808 & - & 0.822 & - \\
\cdashline{2-9}
 & \textit{TVDiag} (released code)$^\ddagger$ & 1.3E-01 & 0.760 & 10.055\% & 0.726 & 11.350\% & 0.742 & 10.767\% \\
 & \textbf{AdaDiag++}$^\dagger$ & 2.1E-04 & \textbf{0.872} & $-$4.144\% & \textbf{0.851} & $-$5.084\% & \textbf{0.862} & $-$4.594\% \\
\bottomrule
\end{tabular}
\par\vspace{4pt}
\begin{minipage}{\linewidth}\footnotesize
$\Uparrow$ indicates the improvement rate of \textit{TVDiag} compared to other methods, as in the paper (negative: the row is above \textit{TVDiag}). The best result in each metric is shown in bold. Rows without $\dagger$ are copied from the paper (Xie et al., \textit{TVDiag}, ACM TOSEM 35(2), Art.~40, 2026, Table 5); the single-modality baselines are omitted. $^\dagger$\textbf{AdaDiag++}: three AdaDiag+ models that differ only in their seed, outputs averaged. AdaDiag+ is TVDiag with adaptive modality weights, each instance's alert-type counts next to its embeddings, and Eq.~13's learnable loss weights (README), mean over 5 seeds (5--9). $^\ddagger$\textit{TVDiag} (released code): the authors' code (github.com/WHU-AISE/TVDiag, commit cf6f334) run unchanged on a CPU, same seeds . p-values of both: two-sided one-sample $t$-test of the seeds against \textit{TVDiag}'s published F1-score.
\end{minipage}
\end{table}

\end{document}
